\begin{itemize}
    \item \textbf{Requisito de Auditoría Física (Restricción Normativa):} El sistema no elimina por completo el uso de papel en la fase de asistencia diaria, debido a la exigencia reglamentaria de los organismos auditores que requieren carpetas físicas con fotocopia de carnet y firma manuscrita para validar la entrega del alimento a cada becado. El software actuará como un medio de control, validación y archivo digitalizador complementario.

    \item \textbf{Operación Interna de la Cocina:} El sistema no gestiona el control de inventarios de insumos, compras de materia prima ni la contabilidad interna de la empresa privada concesionaria del comedor.

    \item \textbf{Comercialización a No Becados:} La venta de almuerzos y platos a estudiantes regulares no becados y personal universitario queda fuera del alcance del sistema, correspondiendo a la gestión propia de la empresa privada.

    \item \textbf{Dependencia de Conectividad:} La carga de expedientes y la sincronización de asistencias requieren conexión a Internet activa.

    \item \textbf{Límites de Verificación Automática:} El sistema realiza una precalificación algorítmica basada en la información declarada y digitalizada; sin embargo, la verificación presencial o visitas domiciliarias mantendrán el criterio operativo del personal de Trabajo Social.
\end{itemize}
